\chapter{最低数学工具与统一记号}
\label{chap:math-tools}

\begin{prereqbox}
本章只整理后文反复使用的最低工具。第一次阅读不必追求形式上的完备；遇到具体章节时，再回到相应小节核对约定和方法。
\end{prereqbox}

\begin{mapbox}
\[
\begin{aligned}
\text{Fourier 变换}&\longrightarrow
\text{倒格子、波动方程、线性响应},\\
\text{本征值与二次型}&\longrightarrow
\text{正规模、能带、平均场},\\
\text{变分与微扰}&\longrightarrow
\text{近似能谱、响应与有效模型},\\
\text{张量和指标}&\longrightarrow
\text{刚体、电磁场与相对论形式}.
\end{aligned}
\]
\end{mapbox}

\section{Fourier 变换与卷积}
统一采用
\[
f(\vb r)=\int\frac{\dd^d q}{(2\pi)^d}
f(\vb q)e^{i\vb q\cdot\vb r},
\qquad
f(\vb q)=\int\dd^dr\,f(\vb r)e^{-i\vb q\cdot\vb r}.
\]
于是
\[
\nabla\longrightarrow i\vb q,\qquad
\partial_t\longrightarrow-i\omega.
\]
平移不变的微分算符在 Fourier 空间中变成代数函数；卷积
\[
(f*g)(\vb r)=\int f(\vb r-\vb r')g(\vb r')\dd^dr'
\]
变成乘积
\[
(f*g)(\vb q)=f(\vb q)g(\vb q).
\]
晶体中的倒格矢展开、Green 函数、衍射和线性响应都使用这一结构。

\section{线性代数、二次型与本征值}
Hermitian 算符满足
\[
A^\dagger=A,
\]
本征值为实数，不同本征值的本征矢正交。若存在简并，本征基在简并子空间中并不唯一；微扰论要求在该子空间内选择能对角化微扰的基。

小振动中的
\[
T=\frac12\dot{\vb q}^{T}M\dot{\vb q},
\qquad
V=\frac12\vb q^{T}K\vb q
\]
给出广义本征值问题
\[
K\vb e_\lambda=\omega_\lambda^2M\vb e_\lambda.
\]
若 $M$ 正定，令
\[
\vb x=M^{1/2}\vb q
\]
即可化为普通 Hermitian 本征值问题
\[
M^{-1/2}KM^{-1/2}\vb u_\lambda
=\omega_\lambda^2\vb u_\lambda.
\]
“选择正确内积后模式正交”比机械计算行列式更重要。

\section{变分法与微扰论的共同逻辑}
变分法在受限试探空间中最小化
\[
E[\psi]=\frac{\langle\psi|H|\psi\rangle}
{\langle\psi|\psi\rangle},
\]
为基态能量给出上界。微扰论则从已知可解问题出发，按小参数系统展开。二者都依赖基底选择：近简并时应先处理简并子空间，强关联时常应从局域占据态而非自由电子态出发。

\section{无量纲化、量纲与极限}
推导完成后至少检查：
\begin{enumerate}
\item 等式两侧量纲是否一致；
\item 相互作用趋零时是否回到可解极限；
\item 温度、波矢或频率趋零时是否符合守恒律；
\item 对称性要求为零的项是否确实消失。
\end{enumerate}

\begin{examplebox}[title=量纲估算自由电子费米能]
零温自由电子气的长度尺度是平均间距
$\ell\sim n^{-1/3}$，动量尺度
\[
p_F\sim\frac{\hbar}{\ell}\sim\hbar n^{1/3}.
\]
因此
\[
E_F\sim\frac{p_F^2}{m}
\sim\frac{\hbar^2}{m}n^{2/3}.
\]
精确态计数只负责补上无量纲系数
$\frac12(3\pi^2)^{2/3}$。
\end{examplebox}

\section{指标、张量与单位制约定}
三维 Euclidean 指标用 $i,j,k$，重复指标默认求和；四维指标用
$\mu,\nu$。叉乘可写为
\[
(\vb a\times\vb b)_i=\epsilon_{ijk}a_jb_k.
\]
正文电动力学默认 SI 制；只有明确标注的相对论专章使用 Gaussian 几何制并设 $c=1$。不同单位制中的电荷、场和四维势不能只靠删除或补写一个 $c$ 相互转换。

\begin{selftestbox}
\begin{enumerate}
\item 在上述 Fourier 约定下，证明卷积变乘积。
\item 为什么广义本征值问题应使用 $M$ 定义的内积？
\item 简并微扰为什么是“先选基底，再算修正”？
\item 用量纲判断三维声学声子的低温热容幂律。
\item 写出恢复 Gaussian 几何制中 $c$ 时必须重新检查的三个量。
\end{enumerate}
\end{selftestbox}
